\subsection{Uniform bounds on weights and vacancies, and interiority}\label{sec:P02}
We retain the notation of Section~\ref{sec:P01}, in particular
$n\ge32$, $\kappa=1/20$, $M_0=6/5$, $d_0=3\log(4/3)$, and
$T=(M_0+d_0+4/e)/\kappa$.

\paragraph{Factorization with the multiplier retained.}
Equation (P1.3) admits the factorization
\[
 x_{ij}=B R_i C_j K_{ij},\qquad
 K_{ij}=h_{D(i,j)}h_{A(i,j)}e^{-\mu_{ij}}(1-p_{ij})^{-3},
                                                               \tag{P2.1}
\]
with positive factors and geometric normalizations
\[
 \frac1n\sum_i\log R_i=\frac1n\sum_j\log C_j=0.
\]
For example, these factors are obtained by dividing $e^{-\alpha_i}$ and
$e^{-\beta_j}$ by their geometric means and absorbing into $B$
these two scalar factors together with $ne^{-2}$.
Each row meets each diagonal in either family at most once.
Thus (P1.5) and the inequality $(1-p)^{-3}\ge1$ give
\[
 \frac1n\sum_j\log K_{ij}\ge-T.
\]
The same property holds for each column:
\[
 \frac1n\sum_i\log K_{ij}\ge-T.                  \tag{P2.2}
\]
Here we use the global sum of $\mu$: the sum of its values in any
row or column is at most the global sum.

\paragraph{Control of the factors without a pointwise vacancy assumption.}
Jensen's inequality, the geometric normalizations, and the row sums give
\[
 1=B R_i\frac1n\sum_j C_jK_{ij}\ge B R_i e^{-T}.
\]
Applying the same argument to the columns gives
\[
 R_i,C_j\le e^T/B.                               \tag{P2.3}
\]
On the other hand, $\log K_{ij}\le d_0$, since $h_d\le1$,
$\mu_{ij}\ge0$, and $p_{ij}\le1/4$.
We multiply the logarithm of (P2.1) by $p_{ij}$ and sum.
The contributions of $\log R_i$ and $\log C_j$ vanish because the
row and column sums are one and the logarithms of the factors in each family sum
to zero. By the convexity used in Section~\ref{sec:P01},
\[
 n\log B=\sum_{ij}p_{ij}\log(np_{ij})
               -\sum_{ij}p_{ij}\log K_{ij}\ge-d_0n.
\]
Thus $B\ge e^{-d_0}$. From (P2.3) we deduce the bound
\[
 B R_i C_j(1-p_{ij})^{-3}
 \le e^{2T+d_0}/B\le e^{2T+2d_0}=:M.             \tag{P2.4}
\]
In (P2.1) all the factors $h_d,e^{-\mu_{ij}}$ are at most
one. Hence, with the caps still in place,
\[
             0<p_{ij}\le M/n.                    \tag{P2.5}
\]

\paragraph{Uniform vacancy bounds.}
Fix a diagonal $d$ in the first family. By (P2.1) and (P2.4),
\[
 m_d(p)=\frac{h_d}{n}\sum_{(i,j)\in d}
 B R_i C_j h_{A(i,j)}e^{-\mu_{ij}}(1-p_{ij})^{-3}
 \le M h_d\frac{|d|}{n}\le M h_d.
\]
The same argument applies to the second family.
Since $m_d=1-h_d$, it follows that for all diagonals
\[
                 h_d\ge\eta:=\frac1{1+M}.        \tag{P2.6}
\]
This conclusion also holds before the cap is removed.

\paragraph{Removal of the cap and a lower bound.}
For
\[
 n\ge n_0:=\max\{32,\lfloor4M\rfloor+1\},
\]
(P2.5) implies $p_{ij}<1/4$. Complementarity then gives
$\mu_{ij}=0$ at every square.
This establishes hypothesis $(H_\infty)$ of Definition~\ref{def:Hinf} and
the uniform bounds. In fact, bounds (P2.5) and
(P2.6) already hold for every $n\ge32$.
For $n\ge n_0$ we also have a lower bound of order $1/n$.
By $K_{ij}\le e^{d_0}$ and (P2.3), the row equation implies
\[
 R_i=\left(B\frac1n\sum_jC_jK_{ij}\right)^{-1}
       \ge e^{-(T+d_0)};
\]
and similarly $C_j\ge e^{-(T+d_0)}$.
Once $\mu=0$, (P2.6) gives $K_{ij}\ge\eta^2$. Together with
$B\ge e^{-d_0}$, this proves
\[
 \frac{c}{n}\le p_{ij}\le\frac Mn,\qquad
 c:=\eta^2e^{-2T-3d_0}>0,
 \qquad h_d\ge\eta.                              \tag{P2.7}
\]

\paragraph{Size of the constants.}
With $\kappa=1/20$, $M_0=6/5$ and $d_0=3\log(4/3)$ we obtain $T\approx70.7$,
$M=e^{2T+2d_0}\approx1.4\cdot10^{62}$, $n_0\approx5.7\cdot10^{62}$ and
$\eta\approx7\cdot10^{-63}$; the constant in~(P2.7) that bounds $n\pi_c$ from below is
$c=\eta^2e^{-2T-3d_0}\approx1.5\cdot10^{-187}$. These constants suffice for the asymptotic statements,
but they do not describe the boards that can be computed: for the ten sizes in Table~\ref{tab:predictor},
from $n=20$ to $n=240$, the numerical solution of Section~\ref{sec:numerico} gives
$0.45\le n\pi_c\le1.41$ and $h_d\ge0.174$. For $6\le n<n_0$ the interiority of the canonical
maximizer is not proved.
